\section{Technologies retenues}

\subsection{Traitement distribué --- Ray}

\gls{ray} est un framework Python de calcul distribué orienté charges \gls{gpu} et workflows \gls{ia}.
Contrairement à Spark (orienté données structurées), Ray est nativement conçu pour
distribuer des inférences de modèles sur GPU, ce qui en fait le choix naturel pour ce projet.

\subsection{Format de données --- ZARR}

\gls{zarr} est un format de tableau N-dimensionnel orienté chunking.
Il permet de lire une image par tuiles sans la charger entièrement en mémoire,
ce qui est critique pour des images haute résolution.
Il est complémentaire à Ray : chaque worker lit uniquement les tuiles qui lui sont assignées.

\subsection{Stockage objet --- à évaluer}

Deux alternatives à \gls{minio} sont à évaluer dans le cadre de ce TB :

\begin{itemize}
    \item \textbf{RustFS} --- alternative S3-only, licence Apache 2.0, migration simple depuis MinIO.
          Moins mature que MinIO mais potentiellement plus performant.
    \item \textbf{CEPH} --- solution généraliste (bloc, fichier, objet), très mature,
          mais complexe à administrer. Pertinent si un cluster CEPH existe déjà.
\end{itemize}

La décision finale sera documentée dans le rapport.

\subsection{Base de données --- PostgreSQL + PostGIS}

\gls{postgresql} avec l'extension \gls{postgis} permet de stocker les polygones
issus de la segmentation avec leurs coordonnées géospatiales.
L'index spatial GIST permet des requêtes géographiques performantes.

Le schéma de base attendu :
\begin{itemize}
    \item table \texttt{annotations} avec colonne \texttt{geom GEOMETRY(POLYGON, 4326)}
    \item index \texttt{GIST} sur la colonne géométrique
\end{itemize}

\subsection{Orchestration --- Kubernetes}

Le cluster \gls{kubernetes} de la HEIG-VD est l'environnement cible de déploiement.
Le pipeline sera packagé sous forme de manifestes Kubernetes (Jobs Ray, Deployments).

\subsection{Modèle --- SAM3}

\gls{sam3} (Segment Anything Model v3) est le modèle de segmentation retenu.
Il sera utilisé en inférence uniquement (pas de réentraînement).
